\subsection{The exact cutoff and the remaining universal lemma}
\label{sec:exact-cutoff}

The obstruction identifies where a hybrid proof must change representations.
Set
\begin{equation}
 \zeta_0:=\lceil\alpha^{-1/4}\rceil.
 \label{eq:separator-cutoff}
\end{equation}
Before the first stage whose live frontal size exceeds $\zeta_0$, the cumulative
dense Schur algebra is at most
\begin{equation}
 O\!\left(\zeta_0^2\operatorname{cvol}(S)\right)
 =O\!\left(\frac{\operatorname{cvol}(S)}{\sqrt\alpha}\right).
 \label{eq:cutoff-dense-work}
\end{equation}
If response queries and
presentation updates are also stopped when their measured charge reaches
this budget, the algorithm reaches the overflow with an exact safe face
state.  This is a rigorous reduction, but not yet a continuation theorem.

For the precise missing statement, write
\begin{equation}
 \Phi_U:=\min\left\{
   \frac12x^TQx-\ell_\rho^Tx:x\geq0,
   \ \operatorname{supp}(x)\subseteq U
 \right\},
 \qquad
 \Delta(U,B):=\Phi_U-\Phi_{U\cup B}\geq0.
 \label{eq:face-shock}
\end{equation}
For an exact positive gate batch, Schur complementation gives
$\Delta(U,B)=q_B^T\mathcal S_B^{-1}q_B/2$.  It is tempting to require a
separate repair after every event with old-face work proportional to
$\log(1+\Delta/E)$.  The next two results show that this target is false.

\begin{theorem}[Explicit RPPR output obstruction]
\label{thm:shock-output-obstruction}
Fix $\alpha\in(0,1)$.  For every sufficiently large $m$ there is a tree and
$\rho=\Theta(1/m)$ whose exact RPPR gate batches begin with a singleton $B$
and then a batch of size $m$, while
\begin{equation}
 \Delta(\{o\},B)=\Theta(1/m^2).
 \label{eq:tiny-shock-large-output}
\end{equation}
Consequently, with $E=\Theta(1/m)$, no explicit-batch algorithm can report
the subsequent gate event in
\begin{equation}
 O\!\left(
  \operatorname{cvol}(B)
  +\frac{\operatorname{cvol}(\{o\}\cup B)}{\sqrt\alpha}
   \log\!\left(1+\frac{\Delta(\{o\},B)}E\right)
 \right)=O(1)
 \label{eq:false-shock-output-bound}
\end{equation}
work.
\end{theorem}

\begin{proof}
Let $a=(1+\alpha)/2$ and $c=(1-\alpha)/2$.  The seed root $o$ has one leaf
$b$ and $m$ neighbors $c_i$, each of which has its own pendant $z_i$.  Thus
$D:=d_o=m+1$, $d_b=d_{z_i}=1$, and $d_{c_i}=2$.  Put
\begin{equation}
 H:=c+2a=1+a,
 \quad g:=\frac c{\sqrt D},
 \quad \sigma:=a-\frac{g^2}{a},
 \quad \lambda:=\frac{g^2}{a\sigma},
 \quad \rho_0:=\frac c{DH},
 \label{eq:shock-tree-parameters}
\end{equation}
and
\begin{equation}
 q_0:=\frac{\alpha c(1-1/H)}{aD},
 \qquad
 \delta:=\frac{a\lambda q_0}{2\alpha(H+\lambda)},
 \qquad
 \rho:=\rho_0+\delta.
 \label{eq:shock-tree-rho}
\end{equation}
At the exact root solve, the violation of $b$ and the common scaled violation
of the $c_i$ are
\begin{equation}
 q_b=\frac\alpha a\left(\frac cD-\rho\right)>0,
 \qquad
 r_c:=\sqrt2e_{\ell_\rho}(x)_{c_i}
 =\frac\alpha a\left(\frac cD-H\rho\right)<0.
 \label{eq:shock-tree-first-gate}
\end{equation}
Hence the first batch is $\{b\}$.  Eliminating $b$ increases every $r_c$ by
$\lambda q_b$, and the choice of $\delta$ gives
\begin{equation}
 r_c+\lambda q_b=\frac{\lambda q_0}{2}>0.
 \label{eq:shock-tree-second-gate}
\end{equation}
Thus the next batch is $\{c_1,\ldots,c_m\}$.  Since
$\lambda=\Theta(1/D)$ and $q_b=\Theta(1/D)$,
$\Delta(\{o\},\{b\})=q_b^2/(2\sigma)=\Theta(D^{-2})$.
Meanwhile $\operatorname{cvol}(\{b\})=2$ and
$\operatorname{cvol}(\{o,b\})=\Theta(D)$.  Taking $E=D^{-1}$ makes the
right-hand side of \cref{eq:false-shock-output-bound} constant, but writing
the next batch requires $m=\Theta(D)$ word outputs.
\end{proof}

Even a compressed batch handle does not rescue a generic full-sweep repair.

\begin{theorem}[No per-event shock bound for full-sweep repair]
\label{thm:no-full-sweep-shock-repair}
Let $A$ have order at least two, and let an enlarged restricted Hessian be
\begin{equation}
 M=\begin{bmatrix}A&g\\g^T&c\end{bmatrix},
 \qquad A\succeq\alpha I,\quad g\leq0,\quad
 \sigma:=c-g^TA^{-1}g\in(0,c).
 \label{eq:full-sweep-enlargement}
\end{equation}
There are a strictly positive old optimum $x^U$, a feasible point
$y=x^U+e$ of old-face error $E$ with $g^Te=0$, and an arbitrarily small
positive new-coordinate demand $q$, such
that optimizing the new coordinate alone leaves enlarged-face error strictly
larger than $E$.  Hence accelerated gradient, Chebyshev, or CG repairs whose
first nontrivial old-face operation is a full matrix--vector product must pay
$\Omega(\operatorname{cvol}(U))$ even as $\Delta/E\to0$.
\end{theorem}

\begin{proof}
Choose a nonzero $e\in\ker(g^T)$ and scale it so that
$e^TAe/2=E$.  Choose $x^U$ positive enough that $x^U+e\geq0$, set the old
load to $Ax^U$, and choose the new load so its demand at $x^U$ is $q>0$.
The exact enlargement gain is $\Delta=q^2/(2\sigma)$.  The demand at $y$ is
still $q$, so optimizing only the new coordinate decreases the objective by
$q^2/(2c)=(\sigma/c)\Delta$.  The enlarged error therefore remains
\begin{equation}
 E+\left(1-\frac\sigma c\right)\Delta>E.
 \label{eq:full-sweep-residual-error}
\end{equation}
Some old coordinate must change.  The stated full-sweep methods cannot make
that change before one old-face matrix--vector product.  Bounded-degree path
or cycle Stieltjes blocks with one pendant new vertex realize
\cref{eq:full-sweep-enlargement}; taking
$\Delta/E=1/\operatorname{cvol}(U)^2$ makes the proposed logarithmic budget
$O(1)$ while that first sweep costs $\Omega(\operatorname{cvol}(U))$.
\end{proof}

The second missing resource is exact sign information.  If $x^T$ is the
exact point on a safe face and $z$ has objective error at most $E$, then
\begin{equation}
 \|z-x^T\|_{Q_{TT}}\leq\sqrt{2E},
 \qquad
 |e_{\ell_\rho}(z)_v-e_{\ell_\rho}(x^T)_v|<\sqrt{2E}
 \label{eq:gate-margin-from-error}
\end{equation}
for every boundary vertex $v$.  Indeed, Cauchy--Schwarz in the $Q_{TT}$ norm
and the positive scalar Schur complement give
$Q_{vT}Q_{TT}^{-1}Q_{Tv}<Q_{vv}\leq1$.  Thus an approximate point certifies
exact gate signs only under a margin larger than $\sqrt{2E}$, or with a
separately charged exact response oracle.

The correct remaining target must therefore be aggregate, explicitly charge
new output and response queries, and allow small shocks to accumulate before
an old-face sweep.

Conditional on exhaustively charging exact response maintenance, the
numerical part of this target is completely amortized.  To avoid hiding work
in an ``oracle,'' define $\mathcal C_{\rm resp}$ to include exactly one state
initialization or conversion at the cutoff, all preprocessing, every exact
sign query and response update on old or newly admitted coordinates, and
every old-face read or arithmetic operation before the final repair.  The
first adjacency-list scan of a newly admitted vertex and its explicit label
output are excluded and charged separately.  The response state contains no
advice depending on $S^\star$ or on future batches.

\begin{theorem}[Aggregate repair given an exact response oracle]
\label{thm:aggregate-repair-oracle}
Let $E>0$, let $U_0\subsetneq\cdots\subsetneq U_J=:T$ be safe faces,
$B_j:=U_{j+1}\setminus U_j$, and
$\Delta_j:=\Phi_{U_j}-\Phi_{U_{j+1}}$.  Let $y\geq0$ be supported on $U_0$
and satisfy
\begin{equation}
 \frac12y^TQy-\ell_\rho^Ty-\Phi_{U_0}\leq E.
 \label{eq:aggregate-initial-error}
\end{equation}
Suppose an online exact-response oracle reports every $B_j$ and certifies
termination with exhaustive charge $\mathcal C_{\rm resp}$ as defined above,
while each admitted adjacency list is otherwise scanned once.  Suppose also
that a fixed-face routine on $T$ satisfies
$G_k\leq C_A\exp(-c_{\rm ctr}k\sqrt\alpha)G_0$ for absolute constants
$C_A\geq1$ and $c_{\rm ctr}>0$ after $k$ full iterations, each costing
$O(\!\operatorname{cvol}(T))$.  Then the segment, including explicit output
and one final repair to error $E$, costs
\begin{equation}
 O\!\left(
  \sum_{j<J}\operatorname{cvol}(B_j)+\mathcal C_{\rm resp}
  +\frac{\operatorname{cvol}(T)}{\sqrt\alpha}\left[
    1+\log\!\left(1+\frac{\sum_{j<J}\Delta_j}{E}\right)
  \right]
 \right).
 \label{eq:aggregate-oracle-work}
\end{equation}
If $\sum_{j<J}\Delta_j\leq C_0E$, its numerical-repair term is
$O_{C_0}(\operatorname{cvol}(T)/\sqrt\alpha)$.
\end{theorem}

\begin{proof}
Do not numerically repair between response events; zero-extend $y$ as the face
grows.  Its objective value does not change, whereas the restricted optimum
telescopes:
\begin{equation}
 \frac12y^TQy-\ell_\rho^Ty-\Phi_T
 \leq E+\Phi_{U_0}-\Phi_T
 =E+\sum_{j<J}\Delta_j.
 \label{eq:aggregate-gap-telescope}
\end{equation}
One final repair therefore needs at most
\begin{equation}
 \left\lceil
  \frac1{c_{\rm ctr}\sqrt\alpha}
  \log\!\left(C_A\left[1+\frac{\sum_{j<J}\Delta_j}{E}\right]\right)
 \right\rceil
 \label{eq:aggregate-final-iterations}
\end{equation}
full iterations.  Its constant logarithm gives the displayed
$\operatorname{cvol}(T)/\sqrt\alpha$ term.  All earlier response operations
belong to $\mathcal C_{\rm resp}$, while the $1$ in
$\operatorname{cvol}(B_j)$ pays each explicit label and first adjacency
scan.  The final $\operatorname{cvol}(T)/\sqrt\alpha$ term also covers
materializing the final vector and its boundary certificate.
\end{proof}

This proves that repeated numerical repair is not the remaining obstacle.
When exact response certificates are maintained independently, heavy shocks
also bound how often a numerical repair is needed.

\begin{theorem}[Aggregate continuation for heavy-shock groups]
\label{thm:heavy-shock-continuation}
Let $E>0$ and let
$U_0\subsetneq\cdots\subsetneq U_H=:T$ be the materialized safe faces of one
accuracy epoch, with $B_h:=U_{h+1}\setminus U_h$ and
$\Delta_h:=\Phi_{U_h}-\Phi_{U_{h+1}}$.  Suppose the epoch begins with
objective error at most $E$ on $U_0$, and
\begin{equation}
 \sum_{h<H}\Delta_h\leq C_{\rm sh}E,
 \qquad
 \Delta_h\geq\theta E\quad(h<H-1),
 \label{eq:heavy-shock-groups}
\end{equation}
where $C_{\rm sh}$ and $\theta>0$ are fixed.  Suppose an independently
maintained exact-response state reports every group and certifies all
nonmembership signs using
$O(\operatorname{cvol}(U_h)+\operatorname{cvol}(B_h))$ work for group $h$,
including all internal response events, and after each group a fixed-face
routine with the contraction in \cref{thm:aggregate-repair-oracle} restores
error at most $E$.  Then the epoch work is
\begin{equation}
 O\!\left(
  \left(\frac{C_{\rm sh}}\theta+1\right)
  \frac{\operatorname{cvol}(T)}{\sqrt\alpha}
  +\sum_{h<H}\operatorname{cvol}(B_h)
 \right).
 \label{eq:heavy-shock-work}
\end{equation}
\end{theorem}

\begin{proof}
The shock floor and telescoping give
$H-1\leq C_{\rm sh}/\theta$.  Immediately after group $h$, zero-padding the
old point gives error at most $E+\Delta_h$.  The fixed-face contraction restores
error $E$ in $O_{C_A,c_{\rm ctr},C_{\rm sh}}(1/\sqrt\alpha)$ iterations.
Every intermediate face has charged volume at most
$\operatorname{cvol}(T)$, so summing those iterations and the assumed exact
response charge per group proves \cref{eq:heavy-shock-work}.  The new batches
are disjoint.
\end{proof}

For a natural exact batch with demand $q_B$ and Schur complement
$\mathcal S_B$, the bounds
$\alpha I\preceq\mathcal S_B\preceq I$ imply
\begin{equation}
 \frac12\|q_B\|_2^2
 \leq\Delta(U,B)
 \leq\frac1{2\alpha}\|q_B\|_2^2.
 \label{eq:shock-demand-sandwich}
\end{equation}
Thus, once an exact response mechanism supplies the natural batches,
\cref{thm:heavy-shock-continuation}'s repair-frequency hypothesis holds if
every nonfinal batch satisfies $\|q_B\|_2^2\geq2\theta E$.  The remaining
\emph{numerical-repair} difficulty is confined to light batches; exact
response certification remains necessary for both heavy and light batches.

\begin{problem}[Graph-uniform light-shock response]
\label{prob:light-shock-response}
Within an accuracy epoch, maintain the exact gate while leaving the numerical
iterate unchanged across consecutive batches with $\Delta_j<\theta E$.
Either terminate or group them until their cumulative shock reaches
$\theta E$, with exhaustive total response charge
\begin{equation}
 \mathcal C_{\rm resp}
 =\widetilde O\!\left(
  \frac{\operatorname{cvol}(S^\star)}{\sqrt\alpha}
 \right)
 \label{eq:light-shock-response-target}
\end{equation}
over all epochs, in addition to explicit output and one scan of each newly
admitted adjacency list.  No terminal-support or future-batch advice is
allowed.
\end{problem}

\begin{lemma}[Graph-universal RPPR support radius]
\label{lem:aggregate-support-radius}
For a finite graph with one seed, put
$R_\star:=\max_{v\in S^\star}\operatorname{dist}(o,v)$, with $R_\star=0$ for
empty support or support containing only the root.  Then
\begin{equation}
 R_\star
 \leq1+
 \frac{\log_+((1-\alpha)/(\alpha\rho))}
 {\log((1+\sqrt\alpha)/(1-\sqrt\alpha))}
 \leq1+
 \frac{\log_+((1-\alpha)/(\alpha\rho))}{2\sqrt\alpha}.
 \label{eq:aggregate-tree-radius}
\end{equation}
\end{lemma}

\begin{proof}
The spectrum of $Q$ lies in $[\alpha,1]$.  Put
$\lambda_\alpha:=(1-\sqrt\alpha)/(1+\sqrt\alpha)$.  The degree-$k$
Chebyshev approximation can be written explicitly by setting
\begin{equation*}
 q_{k+1}(x):=
 \frac{T_{k+1}((1+\alpha-2x)/(1-\alpha))}
      {T_{k+1}((1+\alpha)/(1-\alpha))},
 \qquad p_k(x):=\frac{1-q_{k+1}(x)}x.
\end{equation*}
Because $q_{k+1}(0)=1$, $p_k$ is a polynomial of degree $k$; the extremal
bound for $T_{k+1}$ gives
\begin{equation}
 \|Q^{-1}-p_k(Q)\|_2
 \leq\frac2\alpha\lambda_\alpha^{k+1}.
 \label{eq:aggregate-chebyshev-inverse}
\end{equation}
If $\operatorname{dist}(o,v)=\ell\geq1$, sparsity makes
$(p_{\ell-1}(Q))_{vo}=0$.  For the unregularized point $x^0:=Q^{-1}b$,
$b_o=\alpha/\sqrt{d_o}$ therefore gives
\begin{equation}
 0\leq x_v^\star\leq x_v^0\leq2\lambda_\alpha^\ell.
 \label{eq:aggregate-green-decay}
\end{equation}
Choose a deepest active $v$ when $R_\star\geq1$.  Its active KKT equality,
nonnegativity, and the fact that every neighbor has depth at least
$R_\star-1$ imply
\begin{equation}
 \alpha\rho\sqrt{d_v}
 <\sum_{u\in\mathcal N(v)\cap S^\star}
   \frac{1-\alpha}{2\sqrt{d_ud_v}}x_u^\star
 \leq(1-\alpha)\sqrt{d_v}\lambda_\alpha^{R_\star-1},
 \label{eq:aggregate-deepest-demand}
\end{equation}
where $\sum_{u\in\mathcal N(v)}d_u^{-1/2}\leq d_v$.  Taking logarithms gives
the first inequality in \cref{eq:aggregate-tree-radius}; the second uses
$\log((1+t)/(1-t))\geq2t$.
\end{proof}

\begin{theorem}[Exact aggregate response solver on a rooted tree]
\label{thm:aggregate-tree-response}
Let the RPPR graph be a finite tree rooted at its single seed, let
$\Delta_\star:=\max(\{d_u:u\in S^\star\}\cup\{0\})$, and define $R_\star$ as
in \cref{lem:aggregate-support-radius}.  Exact lazy subtree-response iterators
explicitly emit every active label, perform no old-face sweep during
discovery, and use one final forest resolution.  Their charges satisfy
\begin{align}
 \mathcal C_{\rm resp}
 &\leq O\!\left(
  \log(2+\Delta_\star)\left[
   \operatorname{cvol}(S^\star)
   +\sum_{v\in S^\star}\operatorname{dist}(o,v)
  \right]
 \right),
 \label{eq:tree-aggregate-response}\\
 \Work
 &\leq O\!\left(
  1+\operatorname{cvol}(S^\star)+\mathcal C_{\rm resp}
 \right)
 =\widetilde O\!\left(
  \frac{1+\operatorname{cvol}(S^\star)}{\sqrt\alpha}
 \right).
 \label{eq:tree-aggregate-work}
\end{align}
Thus arbitrary rooted trees bypass the literal gate in
\cref{prob:light-shock-response} with a condition-free exact
response-generated support solver at the same aggregate budget.  On rooted
spiders, \cref{thm:rooted-spider-kinetic} additionally reproduces the literal
exact gate batches.
\end{theorem}

\begin{proof}
For a nonroot $u$ with parent $p$, put $\beta_{pu}:=-Q_{pu}>0$.  Condition on
the parent value $t$, let $z^{(u)}(t)$ be the exact nonnegative minimizer on
the subtree rooted at $u$, and let
$\mu_u(t):=\beta_{pu}z_u^{(u)}(t)$.  If $\mathcal C(u)$ is the child set, the
inverse response of the scalar root coordinate $z=z_u^{(u)}$ is
\begin{equation}
 H_u(z):=\frac{
  Q_{uu}z-\ell_{\rho,u}-\sum_{w\in\mathcal C(u)}\mu_w(z)
 }{\beta_{pu}}.
 \label{eq:aggregate-tree-inverse}
\end{equation}
On each affine piece, its numerator slope is the scalar Schur complement
after eliminating the active descendants, hence
$H_u'(z)\geq\alpha/\beta_{pu}>0$.  Therefore $H_u$ is invertible on its active
part.  Its breakpoints are contained in the merged child breakpoints plus
$u$'s activation, with ties merged; branching creates no new breakpoints.

A lazy iterator stores the current affine piece and one next event per child.
Advancing a consumed event labeled $v$ through a proper support ancestor $u$
costs $O(\log(1+d_u))$, once on every root-path edge.  Initial child records
and the final unconsumed lookahead records cost
$O(\operatorname{cvol}(S^\star)\log(2+\Delta_\star))$: an inactive child is
represented by its edge and degree without scanning its adjacency list, and
any lookahead propagated upward follows the route of its active parent.
This proves \cref{eq:tree-aggregate-response}; each active adjacency list and
emitted label is charged once.

At the root, if $-\ell_{\rho,o}\geq0$, zero satisfies the root and every
nonseed KKT inequality and is returned in $O(1)$ work.  Otherwise the root
performs the same merge on
$Q_{oo}z-\ell_{\rho,o}-\sum_{w\in\mathcal C(o)}\mu_w(z)$ and stops when its
unique zero precedes the next event.  A label is consumed exactly when its
coordinate is positive there, with equality left inactive.  Thus the
consumed labels are precisely $S^\star$.  One top-down forest solve and one
boundary KKT check recover and certify the exact solution in
$O(\operatorname{cvol}(S^\star))$ work.

For nonempty support,
$\sum_{v\in S^\star}\operatorname{dist}(o,v)
\leq R_\star\operatorname{cvol}(S^\star)$ and
$\log(2+\Delta_\star)\leq\log(2+\operatorname{cvol}(S^\star))$.
Combining these bounds with \cref{lem:aggregate-support-radius} proves
\cref{eq:tree-aggregate-work}; the empty-support case was handled separately.
\end{proof}

The tree proof suggests the correct higher-rank sufficient structure, but it
requires an explicit routing contract.

\begin{lemma}[Hierarchical response-route charge]
\label{lem:hierarchical-response-charge}
Suppose an online, locally constructed exact response hierarchy persists
through an epoch.  Every emitted support label $v$ creates one event that
visits each node of a route $P(v)$ once; tied blocks are charged per emitted
label, and there are no unlabeled breakpoints except one live lookahead per
route.  At a hierarchy node $a$, let $r_a$ be the port dimension, let $N_a$
be the maximum number of live response rows, and suppose one event update or
lookahead query costs
\begin{equation}
 c_a=O\!\left(
  r_a^2+\mathcal T_{r_a}(N_a)+\log(1+\deg_{\mathcal H}(a))
 \right).
 \label{eq:hierarchical-node-charge}
\end{equation}
Let $\mathcal C_{\rm build}$ include all initial factorizations, in particular
any $O(r_a^3)$ port setup; let $\mathcal C_{\rm rekey}$ include only rekeys
not already charged in $c_a$; and let $\mathcal C_{\rm recover}$ include
final recovery and certification.  If strict-zero events remain inactive,
then the exact trace generated by the hierarchy has
\begin{equation}
 \mathcal C_{\rm resp}
 \leq
 \mathcal C_{\rm build}+\mathcal C_{\rm recover}
 +\mathcal C_{\rm rekey}
 +\sum_{v\in S^\star}\sum_{a\in P(v)}c_a.
 \label{eq:hierarchical-response-charge}
\end{equation}
\end{lemma}

\begin{proof}
Charge construction and recovery to their named terms.  Persistent state
prevents rebuilding between events.  The routing contract charges every
successful event update once to $(v,a)$; the one live lookahead on a route is
charged to that route's closest emitted predecessor.  Tied events are charged
per label, and all remaining row changes are precisely the separately counted
rekeys.  Summation gives \cref{eq:hierarchical-response-charge}.
\end{proof}

If the right-hand side of \cref{eq:hierarchical-response-charge} is
$\widetilde O((1+\operatorname{cvol}(S^\star))/\sqrt\alpha)$, the hierarchy
meets the aggregate response budget for the exact safe trace it generates.
It solves the literal \cref{prob:light-shock-response} only when its events are
also proved to reproduce those gate batches.  Tree event streams are scalar,
monotone, route once, and have no rekeys; a stable low-dimensional separator
has hierarchy depth one.  The full-response-rank construction following
\cref{eq:kinetic-response-rank} shows why one flat low-rank state cannot
establish \cref{eq:hierarchical-response-charge} on general cyclic cores.

Two further results delimit what a cyclic hierarchy may do.  First, exact
dense fronts cannot be its graph-uniform representation.

\begin{proposition}[Explicit multifrontal hierarchies fail on expanders]
\label{prop:explicit-multifrontal-expander}
On the bounded-degree expander RPPR family of
\cref{prop:expander-response-obstruction}, where
$\alpha_n=\Theta(1/\log n)$, $\rho_n=\Theta(1/n)$, and $S^\star=V$, every
original-basis multifrontal $LDL^T$ or Cholesky hierarchy that explicitly
stores every numerical nonzero of its frontal matrices or triangular factor
has build or storage charge $\Omega(n^2)$.  This exceeds the intended
$\widetilde O(n\sqrt{\log n})$ work even when the terminal support and the
best offline ordering are supplied.
\end{proposition}

\begin{proof}
Let $X$ be a balanced vertex separator.  If $|X|=\Omega(n)$ there is nothing
to prove.  Otherwise, by grouping components of $G-X$, choose their union
$A$ with $n/3\leq|A|\leq2n/3$.  Every edge leaving $A$ is incident on $X$;
bounded degree and edge expansion therefore force $|X|=\Omega(n)$.  Hence
the graph has treewidth $\Omega(n)$, so every scalar vertex-elimination
ordering in the original coordinate basis creates a filled clique of order
$\Omega(n)$.
The fill is numerical, not only symbolic.  Eliminating a Stieltjes pivot
updates
\begin{equation}
 q'_{ij}=q_{ij}-q_{ik}q_{kk}^{-1}q_{kj};
 \label{eq:stieltjes-strict-fill}
\end{equation}
whenever the fill path exists, both incident off-diagonals are negative, so
the new entry is strictly negative and cannot cancel.  The materialized
front or factor therefore contains $\Omega(n^2)$ stored scalars.  Since the
expander has
$\operatorname{cvol}(V)=\Theta(n)$, the claimed separation follows.
\end{proof}

The same issue appears dynamically before a whole dense front is built.

\begin{proposition}[A light cascade defeats explicit demand arrays]
\label{prop:explicit-demand-light-cascade}
There is a family of sparse cyclic SPD Stieltjes systems with an active core
of size $n$ and $n$ exterior candidates such that an exact boundary gate
activates the candidates in $n$ prescribed singleton batches with
arbitrarily small total shock.  After every event, every remaining exact
candidate demand changes.  Consequently any response implementation that
explicitly stores and updates all current scalar candidate demands performs
$\Omega(n^2)$ arithmetic operations or writes, although graph volume and
explicit output are $\Theta(n)$.
\end{proposition}

\begin{proof}
Let $A$ be an irreducible SPD Stieltjes matrix with
$\operatorname{nnz}(A)=O(n)$ whose graph contains a cycle.  Choose $\tau>0$
and $\gamma>\tau^2/\lambda_{\min}(A)$, and set
\begin{equation}
 M:=\begin{bmatrix}A&-\tau I\\-\tau I&\gamma I\end{bmatrix}.
 \label{eq:cyclic-core-leaf-matrix}
\end{equation}
Then $M$ is positive definite.  After the first $n$ coordinates are active,
the
exterior Schur matrix is
\begin{equation}
 S:=\gamma I-\tau^2A^{-1}.
 \label{eq:dense-leaf-schur}
\end{equation}
Irreducibility gives $A^{-1}>0$, so every off-diagonal of $S$ is strictly
negative.  All later Schur complements remain strictly dense Stieltjes by
\cref{eq:stieltjes-strict-fill}.

Work in the reduced exterior quadratic $z^TSz/2-h^Tz$.  Set
$h_1=\varepsilon_1>0$.  Inductively, suppose the first $i-1$ singleton events
have occurred, and let $z^{(k)}$ be the restricted solution on
$\{1,\ldots,k\}$.  For the as-yet unspecified $h_i$, write
$r_i^{(k)}:=h_i-S_{i,1:k}z^{(k)}$.  Strict Stieltjes monotonicity makes
\begin{equation*}
 \delta_i:=r_i^{(i-1)}-r_i^{(i-2)}>0;
\end{equation*}
the difference is independent of $h_i$.  Choose
$0<\varepsilon_i<\delta_i$ and put
\begin{equation}
 h_i=S_{i,1:i-1}z^{(i-1)}+\varepsilon_i.
 \label{eq:light-cascade-load}
\end{equation}
Then $r_i^{(i-1)}=\varepsilon_i>0$ while
$r_i^{(i-2)}<0$; monotonicity gives $r_i^{(k)}\leq0$ at every earlier stage.
Future choices do not alter earlier restricted solutions.  If $\sigma_i>0$
is the current Schur pivot, this event's shock is
$\varepsilon_i^2/(2\sigma_i)$.  Given any $\delta>0$, additionally choosing
$\varepsilon_i<\sqrt{2\sigma_i\delta/2^i}$ makes the total shock smaller
than $\delta$.

To lift the construction to the full system, choose any $\bar x>0$, set
$\ell_U=A\bar x$, and set $\ell_W=h-\tau\bar x$.  Eliminating the active
core $U$ gives exactly the reduced load $h$; after an exterior vector $z$
activates, the core solution is $\bar x+\tau A^{-1}z>0$.  Thus the core
remains active throughout.  Strict density makes event $i$ change all $n-i$
remaining demands.  An eager array that materializes each post-event demand
in a distinct scalar record must therefore perform
$\sum_{i<n}(n-i)=\Theta(n^2)$ scalar updates.  The block system has only
$O(n)$ nonzeros, and the emitted output has $n$ labels.
\end{proof}

These are representation lower bounds, not unconditional computational
lower bounds.  Lazy semiseparable state may compress the cyclic example, and
matrix-free or approximate state may bypass the expander front.  A positive
cyclic composition theorem is nevertheless available when its block
presentation is stable.

\begin{definition}[Stable exposed block--cut presentation]
\label{def:stable-block-presentation}
An online block presentation assigns every exposed edge its true
biconnected-block identifier and every exposed articulation incidence its
parent or child role in the true block--cut tree of the seed component, with
each bridge treated as a size-two block.  The tree is rooted at the seed
vertex when it is an articulation, and otherwise at its unique incident root
block.  Every nonroot vertex label has a unique owner block closest to the
root; a child block conditions on and does not own its parent articulation.
Identifiers, ownership, and roles never merge or change.  Exposing an edge
also supplies its exact normalized coefficient and the neighbor degree.  The
presentation reveals no activation time, terminal-support membership, or
unexposed adjacency.  Let $\mathcal C_{\mathcal B}$ include all work for
revealing and rooting this combinatorial metadata and every presentation-
induced rekey, but not numerical block factors or reverse-solve records.
\end{definition}

\begin{theorem}[Aggregate response with bounded stable blocks]
\label{thm:stable-bounded-block-response}
Suppose a single-seed RPPR graph has a stable exposed block--cut presentation
and the true size of every block incident on an exposed edge is at most $b$.
Put
\begin{equation*}
 N:=1+\operatorname{cvol}(S^\star),
 \qquad
 h_{\mathcal B}(v):=\text{number of block nodes on the root route to $v$}.
\end{equation*}
There is an exact response-generated solver that is adjacency-local relative
to the presentation, with
\begin{equation}
 \Work=O\!\left(
  \mathcal C_{\mathcal B}+b^3N
  +(b^2+\log(2+N))
    \sum_{v\in S^\star}(1+h_{\mathcal B}(v))
 \right).
 \label{eq:stable-bounded-block-work}
\end{equation}
It scans each admitted adjacency list once, emits every support label, charges
all lookaheads, and finishes with exact block-tree recovery and a boundary
KKT certificate.  Since
$h_{\mathcal B}(v)\leq1+\operatorname{dist}(o,v)$, if
$b=\log^{O(1)}(2+N)$ and
$\mathcal C_{\mathcal B}=\widetilde O(N/\sqrt\alpha)$, then
$\Work=\widetilde O(N/\sqrt\alpha)$.
\end{theorem}

\begin{proof}
A directed block--cut subtree $H$ meets its parent only at one articulation
$p$.  Condition on $x_p=t$.  Since $-Q_{Hp}\geq0$, Stieltjes comparison makes
the obstacle solution on $H$ coordinatewise nondecreasing in $t$.  On every
fixed active set it is affine with nonnegative slope, and the scalar response
$-Q_{pH}z(t)$ is continuous, monotone, and piecewise affine.  Every
nontrivial breakpoint carries a newly positive vertex label; ties are grouped
and equality stays inactive.

At a block $B$, current child pieces modify only diagonal and load entries at
their articulation vertices.  These are exact Schur-complement updates, so
the local matrix remains SPD Stieltjes and every inverse-update denominator
is positive.  The active block system has order at most $b$.  A bordered
vertex insertion or one labeled child-piece change updates its inverse in
$O(b^2)$ work, with ties charged per label.  For each articulation, keep a
heap of child breakpoints in the local articulation coordinate and expose
only its minimum to the block heap.  Advancing a child changes one record in
that heap at $O(\log(2+N))$ cost.  The block heap has at most one aggregate
child record per articulation and one activation record per exposed inactive
block vertex, hence $O(b)$ records.  After an inverse update, re-pulling and
rekeying those records costs
$O(b\log b)\subseteq O(b^2+\log(2+N))$.  Exact affine pullback gives the next
parent-parameter event, so each support label advances once through each
ancestor block, with no unlabeled breakpoint proliferation.  A zero-slope
coordinate keeps an unreachable child event at infinity until a later rekey;
immediate and tied events are exhausted before the iterator advances.

At the root, merge the incident block responses and define
\begin{equation*}
 F_o(t):=Q_{oo}t-(\ell_\rho)_o-
         \sum_{B\ni o}\mu_B(t).
\end{equation*}
On every affine piece, $F_o'$ is the scalar Schur complement after eliminating
the active descendants and is at least $\alpha$.  If $F_o(0)\geq0$, zero is
the optimum.  Otherwise consume the next labeled response event whenever it
occurs strictly before the unique positive zero of the current root piece.
Stop when that zero is no later than the next event; equality remains
inactive.  Obstacle comparison then identifies the root (when
$F_o(0)<0$), together with the consumed labels, with exactly $S^\star$ and
gives the exact stopping certificate.

The stable presentation prevents block merging and response reparenting.
Local initialization and root-only or empty-support records are charged to
$N$.  A final unconsumed lookahead propagated through several blocks is
charged to the closest preceding emitted label on the same route; if its
stream has no emitted descendant, it is charged to the active articulation
that owns the stream.  This adds at most a constant copy of the displayed
route sum.  Conservative dense block initialization costs $O(b^3N)$.
Retained bordered and diagonal update records permit reverse block
substitution; one cached scan of active-to-boundary incidences then evaluates
all exact KKT inequalities.  Recovery and certification also cost
$O(b^3N)$.  Therefore the event ledger gives
\cref{eq:stable-bounded-block-work}.  Finally,
\begin{equation*}
 \sum_{v\in S^\star}(1+h_{\mathcal B}(v))
 \leq (2+R_\star)|S^\star|\leq(2+R_\star)N,
\end{equation*}
so \cref{lem:aggregate-support-radius} proves the final claim.  This solver
generates its own safe response trace; no equivalence with the literal
boundary-gate batches is asserted.
\end{proof}

The presentation condition can now be removed for bounded blocks.  The
zero-prefix theorem in \texttt{delayed\_reflection\_ladder} observes that a
newly relevant cycle always contains a just-activated zero coordinate.  The
old response prefix is therefore unchanged; versioned iterators discard only
unconsumed lookaheads, preserve off-path child heaps through articulation
handles, and rebuild the merged at-most-$b$ front from the current checkpoint.
Charging every such rebuild and propagated rekey to the activating incidence
gives ordinary adjacency access the same
$\widetilde O(b^3/(\rho\sqrt\alpha))$ product scale.  Stable metadata remains
a useful abstraction for \cref{thm:stable-bounded-block-response}, but it is
not an oracle assumption for the bounded-block result.

For approximate PageRank output, however, exact zero signs are stronger than
necessary.  The following theorem replaces them by a finite normalized KKT
band.

\begin{theorem}[Finite-resolution aggregate RPPR continuation]
\label{thm:finite-resolution-continuation}
Fix $\rho,\tau_{\rm gate},\tau_{\rm sol}>0$ and put
$\eta:=\alpha\tau_{\rm gate}$.  In the single-seed problem, if
$(\ell_\rho)_o\leq0$, set $T=\varnothing$ and return zero in $O(1)$ work.
Otherwise start
$U_0=\{o\}$ and maintain a \emph{positive Dirichlet face}
\begin{equation*}
 x^U_{U^c}=0,
 \qquad Q_{UU}x^U_U=(\ell_\rho)_U,
 \qquad x^U_U>0.
\end{equation*}
Suppose an online certified response procedure returns, for every
$v\in\partial U$, an estimate of the exact mathematical demand satisfying
\begin{equation}
 |\widehat e_v-e_{\ell_\rho}(x^U)_v|
 \leq\frac\eta4\sqrt{d_v}.
 \label{eq:finite-response-certificate}
\end{equation}
It admits every vertex with
$\widehat e_v>3\eta\sqrt{d_v}/4$, scans each admitted list once, and stops
when no such vertex remains.  Let
$\mathcal C_{\rm resp}^{(\eta)}$ exhaustively charge its preprocessing,
certified error bounds, old-face reads, rekeys, response updates, and final
termination query, excluding the separately charged first scans and label
output.  If $T$ is the terminal face, then
\begin{align}
 T&\subseteq S^\star(\rho),
 &\operatorname{cvol}(T)&\leq\frac2\rho,
 \label{eq:finite-response-safety}\\
 \|D^{-1/2}(x^\star(\rho)-x^T)\|_\infty
 &\leq\tau_{\rm gate}.
 \label{eq:finite-response-gate-error}
\end{align}
If a final feasible fixed-face accelerated solve returns $z\geq0$ supported
on $T$ with
\begin{equation}
 \Psi_{\ell_\rho}(z)-\Psi_{\ell_\rho}(x^T)
 \leq\frac{\alpha\tau_{\rm sol}^2}{2},
 \label{eq:finite-response-solve-error}
\end{equation}
then
\begin{equation}
 \|D^{-1/2}(x^0-z)\|_\infty
 \leq\rho+\tau_{\rm gate}+\tau_{\rm sol}.
 \label{eq:finite-response-total-error}
\end{equation}
The total work is
\begin{equation}
 O\!\left(
  1+\mathcal C_{\rm resp}^{(\eta)}+\operatorname{cvol}(T)
  +\frac{\operatorname{cvol}(T)}{\sqrt\alpha}
    \left[1+\log_+\!\frac{2}{\alpha\tau_{\rm sol}^2}\right]
 \right).
 \label{eq:finite-response-work}
\end{equation}
Every admitted batch additionally has shock density
\begin{equation}
 \Delta(U,B)\geq\frac{\eta^2}{8}\vol(B).
 \label{eq:finite-response-shock-density}
\end{equation}
\end{theorem}

\begin{proof}
If $(\ell_\rho)_o\leq0$, every nonseed load is negative, so zero satisfies
the KKT conditions and is $x^\star$; \cref{eq:rppr-regularization-bias},
proved below, gives the claimed error for the returned zero vector.  Hence
assume $(\ell_\rho)_o>0$.  The initialization is then a positive Dirichlet
face, lies in $S^\star$, and satisfies $0\leq x^{U_0}\leq x^\star$.

Suppose these invariants hold at $U$.  If $v\notin S^\star$, Stieltjes signs
give
$e_{\ell_\rho}(x^U)_v\leq e_{\ell_\rho}(x^\star)_v\leq0$.
Every admitted vertex instead has exact demand larger than
$\eta\sqrt{d_v}/2>0$, so $B\subseteq S^\star$.  Its exact Schur right-hand
side satisfies $q_B>(\eta/2)D_B^{1/2}\one>0$.  Hence
$x_B^{U\cup B}=\mathcal S_B^{-1}q_B>0$, while back-substitution and
Stieltjes signs give $x_U^{U\cup B}\geq x_U^U$.  Since
$U\cup B\subseteq S^\star$, comparison with the active optimum equations
also gives $x^{U\cup B}\leq x^\star$.  This closes the induction.  Every
nonterminal batch is nonempty and contained in the finite set
$S^\star\setminus U$, so the process terminates after at most
$|S^\star\setminus U_0|$ batches.  The support-volume theorem and
\cref{eq:charged-volume-equivalence} give the volume bound.

At termination, \cref{eq:finite-response-certificate} gives
$e_{\ell_\rho}(x^T)_v\leq\eta\sqrt{d_v}$ on the boundary.  The residual is
zero on $T$ and nonpositive outside $T\cup\partial T$.  Put
$h:=x^\star-x^T\geq0$ and let
$M:=\|D^{-1/2}h\|_\infty$.  At a maximizing coordinate $i$ with $M>0$,
$x_i^\star>0$ and
\begin{equation}
 \alpha M\sqrt{d_i}
 \leq(Qh)_i=e_{\ell_\rho}(x^T)_i
 \leq\alpha\tau_{\rm gate}\sqrt{d_i}.
 \label{eq:finite-response-maximum-principle}
\end{equation}
The first inequality uses $h_j\leq M\sqrt{d_j}$, nonpositive
off-diagonals, and $QD^{1/2}\one=\alpha D^{1/2}\one$.  This proves
\cref{eq:finite-response-gate-error}.

For $g:=x^0-x^\star\geq0$, the optimum KKT inequalities give
$(Qg)_i\leq\alpha\rho\sqrt{d_i}$ at every coordinate and equality on the
active support.  Applying the same maximum-principle argument at a maximizing
coordinate of $D^{-1/2}g$ gives
\begin{equation}
 \|D^{-1/2}(x^0-x^\star)\|_\infty\leq\rho.
 \label{eq:rppr-regularization-bias}
\end{equation}
Strong convexity and \cref{eq:finite-response-solve-error} give
$\|z-x^T\|_2\leq\tau_{\rm sol}$, which also bounds the normalized infinity
norm because $d_i\geq1$.  The triangle inequality proves
\cref{eq:finite-response-total-error}.

The zero-start restricted objective gap is at most $1/2$: by
\cref{lem:rppr-unit-envelope}, $\|x^T\|_2\leq1$ and $Q\preceq I$.
Standard fixed-face acceleration therefore proves
\cref{eq:finite-response-work}.  Finally, every admitted exact batch demand
satisfies $q_B\geq(\eta/2)D_B^{1/2}\one$.  Since its Schur complement obeys
$\mathcal S_B\preceq I$,
\begin{equation*}
 \Delta(U,B)=\frac12q_B^T\mathcal S_B^{-1}q_B
 \geq\frac12\|q_B\|_2^2
 \geq\frac{\eta^2}{8}\vol(B),
\end{equation*}
proving \cref{eq:finite-response-shock-density}.
\end{proof}

\begin{corollary}[Margin-adaptive interval continuation]
\label{cor:finite-response-intervals}
In \cref{thm:finite-resolution-continuation}, replace the point estimates by
arbitrary certified intervals
\begin{equation}
 L_v\leq
 \frac{e_{\ell_\rho}(x^U)_v}{\sqrt{d_v}}
 \leq U_v,
 \qquad v\in\partial U.
 \label{eq:finite-response-intervals}
\end{equation}
At a fixed face, admit every label with $L_v>\eta/2$; if none exists, stop
when every $U_v\leq\eta$, and otherwise refine only intervals preventing one
of these two certificates.  If refinement terminates at each fixed face,
then the safety, gate-error, and shock conclusions
\cref{eq:finite-response-safety,eq:finite-response-gate-error,eq:finite-response-shock-density}
all remain valid.  No common interval width is required away from the
transition band.
\end{corollary}

\begin{proof}
The proof of \cref{thm:finite-resolution-continuation} uses its uniform
estimate only twice.  Admission needs the exact inequality
$e_{\ell_\rho}(x^U)_v>\eta\sqrt{d_v}/2$, which follows from the certified
lower endpoint.  Termination needs
$e_{\ell_\rho}(x^T)_v\leq\eta\sqrt{d_v}$, which follows from the certified
upper endpoint.  The same comparison induction, maximum principle, and Schur
shock calculation therefore apply verbatim.
\end{proof}
